\documentclass{article}
\usepackage{amsmath, amssymb, amsthm}
\usepackage{hyperref}
\usepackage{geometry}
\geometry{margin=1in}

\title{Erd\H{o}s Problem \#711}
\date{Last updated: 2025-08-31}
\author{Source: erdosproblems.com}

\begin{document}
\maketitle

\noindent\textbf{Status:} OPEN \quad \textbf{Prize: ₹1000}

\noindent\textbf{Tags:} number theory

\medskip
\noindent\textbf{Problem Statement:}

Let $f(n,m)$ be minimal such that in $(m,m+f(n,m))$ there exist distinct integers $a_1,\ldots,a_n$ such that $k\mid a_k$ for all $1\leq k\leq n$. Prove that\[\max_m f(n,m) \leq n^{1+o(1)}\]and that\[\max_m (f(n,m)-f(n,n))\to \infty.\]

\bigskip
\noindent\textbf{Remarks:}

A problem of Erd\H{o}s and Pomerance \cite{ErPo80}, who proved that\[\max_m f(n,m) \ll n^{3/2}\]and\[n\left(\frac{\log n}{\log\log n}\right)^{1/2} \ll f(n,n)\ll n(\log n)^{1/2}.\]In \cite{Er92c} Erd\H{o}s offered 1000 rupees for a proof of either; for uniform comparison across prizes I have converted this using the 1992 exchange rates.

van Doorn \cite{vD26} has provided an affirmative answer to the second question, proving that, for all large $n$, there exists $m=m(n)$ such that\[f(n,m)-f(n,n) \gg \frac{\log n}{\log\log n}n.\]See also [710].

\bigskip
\noindent\textbf{References:}

[Er92c] Erd\H{o}s, P., Some of my forgotten problems in number theory. Hardy-Ramanujan J. (1992), 34-50.
 
 [ErPo80] P. Erd\H{o}s and C. Pomerance, Matching the natural numbers up to $n$ with distinct multiples of another interval. Indigationes Math. (1980), 147-151.
 
 [vD26] W. van Doorn, On the length of an interval that contains distinct multiples of the first $n$ positive integers. Integers (2026), #A7.

\bigskip
\noindent\small{Source: \url{https://www.erdosproblems.com/711}}
\end{document}
